% Editable standalone vector figure extracted from novogaia_spectra.tex.
% Source slide 2; figure 01; native embedded fonts.
% Compile with LuaLaTeX from this directory, or run ../export_figures.py.
\documentclass[tikz,border=0pt]{standalone}
% Shared editable preamble for standalone, Nature-style vector figures.
% Text remains embedded and searchable. This is a house-style target, not a
% claim of acceptance or of satisfying every journal's submission checklist.
\usepackage{fontspec}
\setmainfont{LiberationSans}[
  Path=../assets/,Extension=.ttf,
  UprightFont=*-Regular,BoldFont=*-Bold,ItalicFont=*-Italic,BoldItalicFont=*-BoldItalic]
\setsansfont{LiberationSans}[
  Path=../assets/,Extension=.ttf,
  UprightFont=*-Regular,BoldFont=*-Bold,ItalicFont=*-Italic,BoldItalicFont=*-BoldItalic]
\setmonofont{LiberationSans}[
  Path=../assets/,Extension=.ttf,
  UprightFont=*-Regular,BoldFont=*-Bold,ItalicFont=*-Italic,BoldItalicFont=*-BoldItalic]
\usepackage{unicode-math}
\setmathfont{Latin Modern Math}
\DeclareMathSizes{5.9}{6}{5}{5}
\DeclareMathSizes{6}{6}{5}{5}
\usepackage{tikz}
\usetikzlibrary{arrows.meta,calc,positioning,patterns,decorations.pathreplacing,matrix}
\usepackage{graphicx}
\usepackage{xcolor}
\pagecolor{white}
\definecolor{ink}{HTML}{202020}
\definecolor{paper}{HTML}{FFFFFF}
\definecolor{muted}{HTML}{494949}
\definecolor{line}{HTML}{808080}
\definecolor{soft}{HTML}{F0F0F0}
\hyphenpenalty=10000\exhyphenpenalty=10000
\newcommand{\figurefont}{\sffamily\fontsize{5.9}{7.4}\selectfont}
\newcommand{\figuremathfont}{\sffamily\fontsize{6}{7.5}\selectfont}
\newcommand{\brandmedium}{\bfseries}
\renewcommand{\normalsize}{\figurefont}
\renewcommand{\small}{\figurefont}
\renewcommand{\footnotesize}{\figurefont}
\renewcommand{\scriptsize}{\figurefont}
\renewcommand{\tiny}{\figurefont}
\renewcommand{\large}{\figurefont}
\renewcommand{\Large}{\figurefont}
\renewcommand{\LARGE}{\figurefont}
\renewcommand{\huge}{\figurefont}
\renewcommand{\Huge}{\figurefont}
\tikzset{
  every path/.style={line width=.55pt},
  every node/.style={font=\figurefont,text=ink},
  arr/.style={-{Latex[length=1.5mm,width=1mm]},line width=.55pt,draw=ink},
  thinarr/.style={-{Latex[length=1.3mm,width=.9mm]},line width=.55pt,draw=ink},
  box/.style={draw=ink,line width=.55pt,fill=white,minimum height=.85cm,align=center,inner sep=3pt},
  ghost/.style={draw=muted,dashed,fill=soft},
  tag/.style={font=\figurefont,text=muted},
  equ/.style={font=\figuremathfont},
  small/.style={font=\figurefont},
  body/.style={font=\figurefont},
  tinytext/.style={font=\figurefont}
}
\newcommand{\txt}[4]{\node[anchor=north west,align=left,text width=#3cm] at (#1,#2) {#4};}
\newcommand{\smalltxt}[4]{\node[anchor=north west,align=left,text width=#3cm,small] at (#1,#2) {#4};}
\newcommand{\eqn}[4][6]{\node[anchor=north west,font=\figuremathfont] at (#2,#3) {$\displaystyle #4$};}
\newcommand{\diagramblock}[5]{\node[box,minimum width=#3cm,text width=#3cm,minimum height=#4cm] at (#1,#2) {#5};}
% Scale coordinates, not text: all labels retain one final publication size.
\newcommand{\spectrum}[3]{\begin{scope}[shift={(#1,#2)}]
  \draw[muted] (0,0)--({3.6*#3},0);
  \draw[muted] (0,0)--(0,{1.9*#3});
  \foreach \x/\h in {.25/.18,.49/.46,.68/1.42,.87/.37,1.08/.74,1.31/.31,1.57/1.78,1.76/.46,1.96/.19,2.17/1.1,2.37/.57,2.57/.31,2.88/.76,3.21/.23,3.4/.38}{
    \draw[ink] ({\x*#3},0)--({\x*#3},{\h*#3});}
  \foreach \x in {.6,1.2,1.8,2.4,3.0}{\draw[ink] ({\x*#3},0)--+(0,{-.06*#3});}
  \foreach \y in {.45,.9,1.35,1.8}{\draw[ink] (0,{\y*#3})--+({-.06*#3},0);}
  \foreach \x/\lab in {0/0,1.8/45,3.6/90}{\node[anchor=north,font=\figurefont] at ({\x*#3},{-.08*#3}) {\lab};}
  \foreach \y/\lab in {0/0,.9/0.5,1.8/1.0}{\node[anchor=east,font=\figurefont] at ({-.08*#3},{\y*#3}) {\lab};}
  \node[anchor=north east,font=\figurefont] at ({3.6*#3},{-.45*#3}) {$m/z$ (Th)};
  \node[anchor=south west,font=\figurefont] at (0,{1.96*#3}) {Relative intensity (a.u.)};
\end{scope}}
\newcommand{\molecule}[3]{\begin{scope}[shift={(#1,#2)}]
  \draw[ink] (0,{#3})--({.866*#3},{.5*#3})--({.866*#3},{-.5*#3})--(0,{-#3})--({-.866*#3},{-.5*#3})--({-.866*#3},{.5*#3})--cycle;
  \draw[ink] ({.68*#3},{.4*#3})--({.68*#3},{-.4*#3});
  \draw[ink] ({-.7*#3},{-.44*#3})--({-.02*#3},{-.83*#3});
  \draw[ink] ({-.7*#3},{.44*#3})--({-.02*#3},{.83*#3});
  \draw[ink] (0,{#3})--(0,{1.55*#3});
  \node[font=\figurefont] at (0,{1.77*#3}) {OH};
\end{scope}}
\newcommand{\peaktokens}[3]{\foreach \i in {0,...,5}{\pgfmathsetmacro{\xx}{#1+\i*.36}\draw[fill=#3,draw=ink] (\xx,#2) rectangle +(.24,.6);}}
\newcommand{\E}{\mathbb E}
\newcommand{\R}{\mathbb R}
\newcommand{\sg}{\operatorname{stopgrad}}
\newcommand{\LN}{\operatorname{LN}}
\newcommand{\softmax}{\operatorname{softmax}}

\begin{document}
\figurefont
\begin{tikzpicture}[x=1cm,y=1cm,scale=1.14375,transform shape]
\molecule{2.05}{4.95}{.7}
\node[tag] at (2.05,3.58) {$G=(V,A)$};
\draw[arr] (3.24,5.54)--node[above,small]{fragmentation}(6.1,5.54);
\draw[arr,dashed] (6.1,4.56)--node[below,small]{inverse problem}(3.24,4.56);
\spectrum{6.7}{4.13}{.9}
\node[tag] at (8.32,2.98) {$S=\{(m_i,I_i)\}_{i=1}^{N}$};
\node[small,anchor=base,inner sep=0] at (13.25,6.42) {Same formula: $F=\mathrm{C_6H_6O}$};
\coordinate (candidatefork) at (11.1,4.75);
\draw[arr] (10.32,4.75)--(candidatefork);
\draw[ink,line width=.55pt] (11.1,3.2)--(11.1,5.72);
\foreach \y in {5.72,4.46,3.2}{\draw[thinarr] (11.1,\y)--(11.91,\y);}
% Phenol, SMILES Oc1ccccc1. Alternating inner bonds retain line-angle clarity.
\begin{scope}[shift={(12.83,5.72)}]
 \draw[ink,line width=.55pt] (-.42,0)--(-.21,.364)--(.21,.364)--(.42,0)--(.21,-.364)--(-.21,-.364)--cycle;
 \draw[ink,line width=.55pt] (-.327,.016)--(-.167,.294) (.167,.294)--(.327,.016) (.17,-.284)--(-.17,-.284);
 \draw[ink,line width=.55pt] (.42,0)--(.73,0);
 \node[anchor=west,inner sep=0,small] at (.79,0) {OH};
\end{scope}
\node[anchor=base west,inner sep=0,tag] at (14.7,5.66) {$G_1$};
% Cyclohexa-2,4-dien-1-one, SMILES O=C1C=CC=CC1.
\begin{scope}[shift={(12.83,4.46)}]
 \draw[ink,line width=.55pt] (-.42,0)--(-.21,.364)--(.21,.364)--(.42,0)--(.21,-.364)--(-.21,-.364)--cycle;
 \draw[ink,line width=.55pt] (-.327,.016)--(-.167,.294) (.17,-.284)--(-.17,-.284);
 \draw[ink,line width=.55pt] (.42,.038)--(.75,.038) (.42,-.038)--(.75,-.038);
 \node[anchor=west,inner sep=0,small] at (.81,0) {O};
\end{scope}
\node[anchor=base west,inner sep=0,tag] at (14.7,4.40) {$G_2$};
% Hexa-2,4,5-trienal, SMILES O=CC=CC=C=C; the cumulated bonds are collinear.
\begin{scope}[shift={(12.65,3.12)}]
 \node[anchor=south east,inner sep=0,small] at (-.265,.157) {O};
 \draw[ink,line width=.55pt] (-.234,.179)--(.019,.031) (-.272,.113)--(-.019,-.031);
 \draw[ink,line width=.55pt] (0,0)--(.34,.20)--(.68,0)--(1.02,.20)--(1.36,0)--(1.70,-.20);
 \draw[ink,line width=.55pt] (.381,.141)--(.682,-.035) (1.061,.141)--(1.741,-.259);
 \node[fill=paper,inner sep=.65pt,font=\figurefont] at (1.36,0) {C};
\end{scope}
\node[anchor=base west,inner sep=0,tag] at (14.7,3.14) {$G_3$};
\node[anchor=west,inner sep=0,font=\figurefont] at (1.03,2.63) {$\displaystyle\widehat G=\arg\max\nolimits_G p(G\mid S,F)$};
% Fixed canvas: 183 x 60.046875 mm. No clipping: overflow is audited.
\pgfresetboundingbox
\path[use as bounding box] (0,1.6) rectangle (16,6.85);
\end{tikzpicture}
\end{document}
